\documentclass[12pt, a4paper, oneside]{ctexart}
\usepackage{amsmath,amsthm,amssymb,graphicx}
\usepackage{float}
\usepackage[margin=2.5cm]{geometry}
\usepackage{xcolor}
\usepackage{enumitem}

\begin{document}
我们在之前的课程中讲过的加解密过程是十分自然的：一个人拿着自己的密钥，想加密就用这个密钥加密，想解密就用这个密钥解密，看上去这是天衣无缝的，但是在某些情况下，这并不安全，比如说Alice掌管着公司里的一份密钥，如果Alice被敌对公司策反了，那么敌对公司就可以使用Alice手中的密钥获得公司内部的一切信息，仅仅因为Alice一人的背叛就威胁到了一整个公司的信息安全，这显然是不合理的，所以针对这种情况，我们需要一种新的加密方式，使得仅凭一个人的力量并不能实现加解密。

那我们很自然地就可以想到，那我们直接把一份密钥拆成好几份，分发给不同的人不就OK了吗？这样仅仅一个人的叛变也无法造成信息的泄露，但是这也有一个问题，那就是如果我们把密钥拆成了好几份，那么我们就需要所有人都到齐才能加解密，这显然是不合理的，所以我们需要一种新的加密方式，使得仅凭一部分人的力量就可以实现加解密。


这就是我们下面要提到的$(t,n)$秘密共享方案：分发者$D$将秘密$s$拆分成$n$个份额$s_1,s_2,\ldots,s_n$，分别分发给$n$个参与者$P_1,P_2,\ldots,P_n$，每个参与者$P_i$接收到一个份额$s_i$，其中$i\in [1,n]$,使得任意$t$个参与者就可以联合起来恢复出秘密$s$，而任意少于$t$个参与者合作不能获得任何关于$s$的信息。

乍一看这个方案好像很容易实现，比如说实现(2,2)秘密共享：我要是有一个128位的密钥，我把前64位分给Alice，把后64位分给Bob不就实现了吗？看上去合理实际上一点也不，我们的方案里面写过，“任意少于$t$个参与者合作不能获得任何关于$s$的信息”，所以现在这样分配显然是不行的，因为不论是Alice还是Bob都获得了密钥的部分信息，这样做最直接的效果就是让密钥穷举搜索的空间变小了，叛变的一方可以直接穷举另一方的$2^50$种可能来获得完整的密钥，这显然是不安全的。我们来看看下面这个方案：

还是一个128位的密钥$s$，我们先随机选择一个128位的随机数$s_1$，并令$s_2=s\oplus s_1$，我们将$s_1$给Alice，把$s_2$给Bob，这样就很合理了，因为不论是Alice还是Bob，都无法从自己的份额中拿到关于密钥的信息（被随机数掩盖掉了），也没有办法穷举（因为此时的穷举空间是$2^{128}$），两人合在一起恰好能恢复密钥$s=s_1\oplus s_2$。


























\end{document}